%! AP Statistics — Unit 8, Lesson 8.1 — Student Packet Cover Sheet
\documentclass[10pt]{article}
\usepackage{apstats-article}
\usepackage{apstats-boxes}
\usepackage{tikz}

\begin{document}

\begin{tikzpicture}[remember picture, overlay]
  \fill[navy] (current page.north west) rectangle
    ([yshift=-2.2cm]current page.north east);
  \node[anchor=west, white, font=\large\bfseries] at
    ([xshift=1cm, yshift=-0.7cm]current page.north west)
    {\CourseName: \SchoolYear};
  \node[anchor=west, white, font=\normalsize] at
    ([xshift=1cm, yshift=-1.3cm]current page.north west)
    {Unit 8: Inference for Categorical Data: Chi-Square};
  \node[anchor=west, white, font=\normalsize\bfseries] at
    ([xshift=1cm, yshift=-1.9cm]current page.north west)
    {Lesson 8.1 \quad Introducing Statistics: Are My Results Unexpected?};
\end{tikzpicture}

\vspace{2.4cm}
\namedateperiod

\begin{learningtargetbox}
\begin{itemize}
    \item I can identify a question about whether variation between observed and
          expected counts is due to chance.
    \item I can explain what it means for categorical data to follow an expected
          distribution.
    \item I can describe how chi-square tests address questions about categorical
          data distributions.
\end{itemize}
\end{learningtargetbox}

\begin{tocbox}
% A CONTENTS LIST ONLY: no Score column and no Total row. Whether a component counts is said in
% its description.
\renewcommand{\arraystretch}{1.5}
\begin{tabularx}{\linewidth}{clX}
    \textbf{\#} & \textbf{Component} & \textbf{Description} \\
    \hline
    1 & Warm-Up        & Spiral review: categorical distributions, proportions, and expected counts \\
    2 & Guided Notes   & Vocabulary, expected counts, the chi-square idea, unit arc overview \\
    3 & Activity       & Group work: observed vs.\ expected counts in a survey scenario \\
    4 & Exit Ticket    & Independent check on expected counts and interpreting differences \\
    5 & Homework       & Practice with expected counts and chi-square framing \\
\end{tabularx}
\end{tocbox}

\end{document}
